La validación en placa comprueba el funcionamiento del sistema en las
condiciones reales de uso: la PC envía programas y órdenes por UART, y la FPGA
ejecuta las instrucciones y devuelve su estado. Esto complementa las pruebas de
simulación presentadas anteriormente y permite observar el recorrido completo,
desde el código assembler hasta las respuestas recibidas desde la placa.

\section{Configuración de la validación}

Para la validación se conectó una PC a la Basys 3 y se utilizó el assembler
junto con la aplicación \tech{viewrv}. La FPGA se configuró con el bitstream
del perfil \code{REFERENCE}; sobre esa misma configuración se cargaron y
ejecutaron diferentes programas, sin volver a sintetizar el diseño.

Se seleccionaron 151 variantes de pruebas ya definidas en el catálogo. La
tabla~\ref{tab:replay} indica a qué grupo pertenece cada conjunto y cuántas
variantes se seleccionaron. Se conservan los identificadores originales para
poder comparar lo observado en placa con las expectativas de cada prueba.

\begin{table}[H]
  \centering
  \small
  \begin{tabular}{llr}
    \toprule
    \textbf{Grupo} & \textbf{Tema} & \textbf{Variantes} \\
    \midrule
    \code{TEST-ISA-001} & Aritmética y lógica & 20 \\
    \code{TEST-ISA-002} & Shifts y comparaciones & 37 \\
    \code{TEST-ISA-003} & \code{LUI} y límites & 4 \\
    \code{TEST-ISA-005} & Branches condicionales & 4 \\
    \code{TEST-ISA-006} & Saltos y valor de enlace & 8 \\
    \code{TEST-ISA-007} & Loads y stores & 22 \\
    \code{TEST-ISA-008} & Destinos de flujo inválidos & 6 \\
    \code{TEST-PIPE-003} & Forwarding y prioridad por fuente & 5 \\
    \code{TEST-PIPE-005} & Consumidores load-use & 9 \\
    \code{TEST-EXEC-004} & Faults precisos & 7 \\
    \code{TEST-EXEC-008} & Ausencia de halt & 4 \\
    \code{TEST-MEM-008} & Reemplazo y publicación de imagen & 6 \\
    \code{TEST-UART-004} & Carga de programas & 9 \\
    \code{TEST-UART-008} & Lista de DMEM & 4 \\
    \code{TEST-UART-013} & Secuencia de reset & 6 \\
    \midrule
    Total & & 151 \\
    \bottomrule
  \end{tabular}
  \caption{Variantes seleccionadas para la validación en placa.}
  \label{tab:replay}
\end{table}

Para preparar los datos de cada prueba se ejecutaron instrucciones reales en la
CPU. Por ejemplo, los casos de \code{TEST-UART-008} utilizan stores para
escribir DMEM y generar la lista de memoria utilizada. Las instrucciones de
preparación forman parte de la ejecución, por lo que sus efectos y sus ciclos se
tienen en cuenta al comparar los resultados.

La selección conserva el alcance previsto para la placa: en
\code{TEST-UART-013}, por ejemplo, los cuatro casos de drenado se verifican
localmente y no forman parte de las variantes físicas seleccionadas.

La tabla~\ref{tab:sesion-fisica} resume la configuración utilizada para la
comunicación y la ejecución en placa.

\begin{table}[H]
  \centering
  \small
  \begin{tabularx}{\textwidth}{L{3.8cm}Y}
    \toprule
    \textbf{Elemento} & \textbf{Configuración} \\
    \midrule
    Placa & Basys 3. \\
    Enlace serie & \num{19200} baud, 8N1, sin control de flujo. \\
    Bitstream & \code{tp3-reference.bit}, perfil \code{REFERENCE}. \\
    Aplicación & \tech{viewrv} 0.7.0. \\
    \bottomrule
  \end{tabularx}
  \caption{Configuración utilizada para la validación en placa.}
  \label{tab:sesion-fisica}
\end{table}

\section{Resultados de la validación en placa}

De las 151 variantes seleccionadas, 150 aprobaron en placa: 146 se ejecutaron
automáticamente y cuatro requirieron accionar el botón BTNC durante RUN, LOAD
o Debug. Estas pruebas permitieron comparar las respuestas y el estado observado
con los resultados esperados para cada caso.

La variante \code{TEST-UART-013/C009} quedó pendiente en placa. Esta prueba
requiere presionar el botón de reset durante la preparación posterior a una
carga, que dura aproximadamente \SI{41}{\micro\second}. Para asegurar que el
pulso ocurra en esa ventana es necesario sincronizarlo con el final de la carga
y registrar su posición temporal. Como no se contó con esa instrumentación, la
variante no se contabiliza entre las aprobaciones físicas.

La misma variante sí aprobó en simulación de sistema con el banco
\code{system_top_session_tb}, que comprobó que la preparación estuviera activa
al llegar el reset, que la transmisión volviera a reposo y que el sistema
quedara en un estado limpio. El resultado corresponde a simulación y no
sustituye la prueba pendiente en placa.

\section{Recorrido funcional comprobado}

Además de las variantes de la tabla, se comprobó un recorrido de uso mediante
la interfaz de línea de comandos de \tech{viewrv}. Este recorrido utilizó el
assembler, la aplicación y la FPGA para cargar programas, ejecutarlos y
consultar su estado.

Para describirlo se identifican como A y B dos programas de prueba. El programa
A escribe valores en registros y memoria, mientras que B termina sin dejar
datos escritos. La secuencia permitió observar el avance de A y después
comprobar que B comenzara con un estado inicial limpio:

\begin{enumerate}
  \item \textbf{Preparación y carga:} se configuró la FPGA, se comprobó el
    enlace y se ensambló y cargó el programa A.
  \item \textbf{Paso a paso:} el primer \code{STEP} produjo ciclo 1, PC 4,
    \code{IF/ID} válido en PC 0, RF en cero y DMEM sin bytes utilizados.
  \item \textbf{Ejecución continua:} tras recargar A, \code{RUN} terminó en
    \code{FINISHED}, ciclo 7, PC 20, \code{x1=7} y \code{x2=9}; los cuatro
    bytes utilizados de DMEM fueron \code{07 00 00 00} en las direcciones 0--3.
  \item \textbf{Reprogramación:} se cargó B sin cambiar el bitstream. Terminó en
    \code{FINISHED}, ciclo 3, PC 8, RF en cero y DMEM sin bytes utilizados.
  \item \textbf{Fault:} un tercer programa terminó en \code{FAULT}, ciclo 5,
    PC 8, causa 1, \code{fault_pc=4} y \code{x1=1}.
  \item \textbf{Interrupción y recuperación:} se inició un bucle sin
    \code{halt}, se solicitó su detención desde otra terminal y se confirmó
    la interrupción por petición del usuario. La carga y ejecución posterior
    de B comprobó la recuperación.
\end{enumerate}

Además, las variantes seleccionadas incluyeron programas con dependencias,
comparaciones de snapshots y pruebas del botón de reset. En particular, los
casos de \code{load-use} permitieron comprobar los ciclos y campos activos
esperados para esa dependencia.

\section{Alcance de la verificación de \texorpdfstring{\tech{viewrv}}{viewrv}}

La aplicación \tech{viewrv} tiene sus propios requisitos y pruebas. Mientras
que las variantes de \code{TEST-PLAT-005} se concentran en el comportamiento de
la FPGA observado por UART, los escenarios de la aplicación comprueban cómo se
preparan las operaciones, se interpretan las respuestas y se presentan los
resultados. El \doc{software/viewrv/verificacion.md}{catálogo de la aplicación}
define 30 escenarios y 74 requisitos.

Se realizaron 47 pruebas de protocolo y de la CLI con el core, mediante
adaptadores controlados y pseudo-terminales de Linux. Estas pruebas permiten
comprobar, sin una placa conectada, la formación de mensajes, las respuestas
parciales, los tiempos de espera y la recuperación de la comunicación. El
recorrido físico descrito anteriormente agrega la observación del flujo con una
Basys 3 real.

La interfaz gráfica comparte el core y presenta el código, los registros, el
pipeline y la memoria. Sus pruebas unitarias y recorridos de ventana se
describen en la \doc{software/viewrv/verificacion.md}{verificación de viewrv}.
La ejecución física mediante
CLI no acredita por sí sola todos los escenarios de la interfaz gráfica ni la
distribución para Windows.

\section{Registro de evidencia}

Para cada variante se conserva el plan de prueba y se registra el resultado
esperado y el observado. Según el caso, la evidencia incluye los bytes
intercambiados por UART y el snapshot recibido, lo que permite revisar los
registros, el pipeline, la memoria utilizada y la información de fault.

El \evidencefile{resultados-por-id.csv}{resumen por ID} reúne los resultados
de las variantes y el \evidencefile{tp3-issue-17-evidencia.zip}{paquete de
evidencia} conserva las capturas, el bitstream y los reportes utilizados.

\section{Resultado de aceptación}

La validación comprobó en placa la carga y el reemplazo de programas, la
ejecución continua y paso a paso, la observación del procesador y la
recuperación mediante reset. Los resultados permiten contrastar la ejecución
física con las expectativas de las pruebas, teniendo en cuenta la variante de
reset que permanece pendiente en placa.
